\section{Round-six positive closure: compound-Poisson smoothing at the exact finite saddle}
\label{sec:r6-b1}

The conditioning change of measure is centered at the exact finite-volume
mean.  Smoothing of momentum and energy is derived from the singleton
compound-Poisson exponent of the grand-canonical logarithm; no collection of
independent particles is inserted as an additional connected cluster.

\subsection{Constraint coordinates and the finite pressure}

Let
\[
C(x,v)=\left(1,v,\frac{|v|^2}{2},\chi_1(x,v),\ldots,\chi_m(x,v)\right)
\in\mathbb R^d.
\]
The first coordinate is lattice valued after multiplication by the
large-deviation speed \(\mu_\varepsilon\); the remaining coordinates are
continuous.  We impose the following explicit nondegeneracy conditions.

\begin{enumerate}[label=(C\arabic*)]
\item the particle-number lattice has span one;
\item after conditioning on number, the pushforward of the singleton activity
by the continuous part of \(C\) has a \(C^{d+4}\) density on one open set and
its affine span is \(\mathbb R^{d-1}\);
\item no nonzero continuous frequency has modulus-one singleton
characteristic function on all positive-activity patches;
\item the target vector lies in a compact subset of the interior mean image.
\end{enumerate}

For a particle-path/contact source \(H\) and multiplier \(\lambda\), put
\[
Q_\varepsilon(H,\lambda)=
\frac1{\mu_\varepsilon}
\log\mathbb E_\varepsilon^{\rm gc}
\exp\left\{
\mu_\varepsilon\langle H,(\Pi^\varepsilon,\Gamma^\varepsilon)\rangle
+\lambda\cdot\sum_iC(z_i(0))
\right\}.
\]
The B2 grand-canonical theorem supplies a normal connected expansion on every
bounded real source ball and on a complex neighborhood of its interior.

\subsection{The exact singleton exponent}

Let \(a_{H,\lambda}^\varepsilon(dz)\) be the one-label activity after the
one-particle part of the dynamic source has been integrated along a free
trajectory.  The connected logarithm separates exactly as
\[
Q_\varepsilon(H,\lambda+i\theta)-Q_\varepsilon(H,\lambda)
=\mathfrak p_{1,\varepsilon}(\theta)
+\mathfrak r_\varepsilon(H,\lambda,\theta),
\]
where
\[
\mathfrak p_{1,\varepsilon}(\theta)=
\int\left(e^{i\theta\cdot C(z)}-1\right)
 a_{H,\lambda}^\varepsilon(dz)
\]
is the compound-Poisson singleton exponent and
\(\mathfrak r_\varepsilon\) is the sum of connected clusters with at least
two labels.

\begin{lemma}[Compound-Poisson smoothing and connected remainder]
\label{lem:r6-b1-poisson}
On compact source and multiplier sets there are constants
\(c,C,\delta>0\) such that:
\begin{enumerate}[label=(\roman*)]
\item the singleton activity has total mass bounded above and below and has
uniform \(C^{d+4}\) densities on the patches in (C2);
\item for \(|\theta|\le\delta\),
\[
\Re\mathfrak p_{1,\varepsilon}(\theta)
\le-c|\theta|^2;
\]
\item on every compact frequency set separated from the exact mixed lattice
origin,
\[
\Re\mathfrak p_{1,\varepsilon}(\theta)\le-c;
\]
\item for continuous frequencies \(|u|\ge1\), the exact compound-Poisson
coefficient obtained from
\[
\exp\{\mu_\varepsilon\mathfrak p_{1,\varepsilon}(t,u)\}
=e^{-\mu_\varepsilon a_1}
\sum_{n\ge0}\frac{\mu_\varepsilon^n}{n!}
\left(\widehat a_{H,\lambda}^\varepsilon(t,u)\right)^n
\]
has an integrable Fourier majorant of arbitrary prescribed polynomial order;
\item the connected remainder and its first four frequency derivatives obey
\[
|\partial_\theta^j\mathfrak r_\varepsilon|
\le C(T+\varepsilon),
\qquad 0\le j\le4.
\]
\end{enumerate}
\end{lemma}

\begin{proof}
The one-label term in a grand-canonical logarithm is exactly the singleton
activity integral.  Exponentiating it gives the displayed Poisson series; this
is where convolution powers of singleton marks occur.  They are not added to
the connected logarithm.  Conditions (C2)--(C3) and compactness give the
quadratic and compact-frequency Cramer bounds for the singleton characteristic
function.  A convolution power of order at least \(d+4\) has an integrable
\(C^4\) density by the coarea rank condition.  Repeated integration by parts
therefore gives arbitrary polynomial Fourier decay for the corresponding
terms of the exact Poisson series.  The finitely many smaller powers are
handled directly and the Poisson tail is exponentially summable.

Every cluster in \(\mathfrak r_\varepsilon\) contains at least one connecting
interaction or contact and hence one factor \(T+\varepsilon\) in the B2
trace-class cluster norm.  Mark insertions differentiate only bounded
constraint coordinates on the declared source chart.  Normal convergence
gives the uniform derivative bound.  Decreasing the common short-time
horizon, the singleton gaps dominate the connected remainder.
\end{proof}

\subsection{Uniform finite-volume convexity and mean image}

\begin{lemma}[Number/mark covariance block]
\label{lem:r6-b1-covariance}
The finite-volume Hessian admits the block decomposition
\[
D^2_{\lambda\lambda}Q_\varepsilon=
\begin{pmatrix}
\sigma_N^2&b^T\\b&\Sigma_Y
\end{pmatrix},
\]
where \(\sigma_N^2\ge c\) comes from activity fluctuations and the Schur
complement
\(\Sigma_Y-bb^T/\sigma_N^2\ge cI\).  Hence
\[
cI\le D^2_{\lambda\lambda}Q_\varepsilon(H,\lambda)
\le CI
\]
uniformly on compact source charts.
\end{lemma}

\begin{proof}
The constant coordinate is not centered inside a fixed-particle singleton
law; its fluctuation is the compound-Poisson activity count.  This gives the
positive number block.  Conditional on the number sector, the continuous
mark covariance is positive by (C2) and the spanning activity patches.  The
Schur complement is therefore positive for the singleton exponent.  The
connected correction to the Hessian is \(O(T+\varepsilon)\) by
Lemma~\ref{lem:r6-b1-poisson}; choose the common horizon so it is smaller than
half the singleton lower bound.  Upper bounds follow from Gaussian velocity
tails and normal convergence.
\end{proof}

Let \(a_\varepsilon\) be the exact finite shell center, including the integer
particle number divided by \(\mu_\varepsilon\).

\begin{theorem}[Exact finite-volume saddle]
\label{thm:r6-b1-saddle}
For every source \(H\) in a compact source ball and every compatible target
\(a_\varepsilon\) in the fixed interior set, there is a unique multiplier
\(\lambda_{H,\varepsilon}\) such that
\[
D_\lambda Q_\varepsilon(H,\lambda_{H,\varepsilon})=a_\varepsilon.
\]
The multipliers remain in one compact interior chart, converge locally
uniformly to the unique limiting saddle \(\lambda_H\), and are analytic in
\(H\).
\end{theorem}

\begin{proof}
The limiting gradient maps the compact multiplier chart diffeomorphically
onto an open set containing the target compact set.  Uniform convergence of
the finite gradients on the chart places the same compact target inside every
finite mean image for small \(\varepsilon\).  The uniform Hessian lower bound
from Lemma~\ref{lem:r6-b1-covariance} gives injectivity and the quantitative
inverse-function theorem gives existence, uniqueness, and analytic source
dependence.  Compactness and uniqueness identify the limit.
\end{proof}

\subsection{Mixed lattice--continuous local limit at the finite saddle}

Under the \((H,\lambda_{H,\varepsilon})\)-tilted law, write
\[
\widehat N=N-N_\varepsilon,
\qquad
\widehat Y=Y-\mathbb E_{H,\lambda_{H,\varepsilon}}Y.
\]
By construction both means are exactly centered at the target shell.

\begin{theorem}[Uniform mixed characteristic-function theorem]
\label{thm:r6-b1-characteristic}
There are \(c,\delta>0\) such that the normalized joint characteristic
function \(\varphi_\varepsilon(t,u)\) satisfies
\[
|\varphi_\varepsilon(t,u)|
\le
\begin{cases}
\exp\{-c\mu_\varepsilon(t^2+|u|^2)\},
&|t|+|u|\le\delta,\\
\exp\{-c\mu_\varepsilon\},
&\delta\le d((t,u),\Lambda^*)\le C,\\
C_M(1+\sqrt{\mu_\varepsilon}|u|)^{-M},
&|u|\ge C,
\end{cases}
\]
where \(t\in[-\pi,\pi]\), \(\Lambda^*\) is the exact mixed dual lattice,
and the fundamental domain contains only the origin.  The estimates and their
first two source derivatives are uniform on compact source and target sets.
\end{theorem}

\begin{proof}
Near zero, Taylor expansion at the exact finite saddle and
Lemma~\ref{lem:r6-b1-covariance} give the Gaussian bound with no mean error.
On compact minor arcs, Lemma~\ref{lem:r6-b1-poisson} gives a strict singleton
gap.  The connected remainder is smaller on the common short horizon.  The
span-one and nonlattice assumptions identify the origin as the only dual
lattice point.  For large continuous frequency, use the exact Poisson series
and its smooth convolution powers from
Lemma~\ref{lem:r6-b1-poisson}; the connected factor has bounded derivatives
and is absorbed by convolution.  Source derivatives insert only integrable
marks, so the same majorants apply.
\end{proof}

Let \(B\subset\mathbb R^{d-1}\) be a fixed bounded convex set with piecewise
\(C^2\) boundary and nonempty interior.  Let
\(\delta_\varepsilon\downarrow0\) satisfy
\[
\sqrt{\mu_\varepsilon}\delta_\varepsilon\to\infty.
\]

\begin{theorem}[Exact-number shrinking-shell coefficient]
\label{thm:r6-b1-coefficient}
Uniformly on compact source and target sets,
\[
\mathbb P_{H,\lambda_{H,\varepsilon},\varepsilon}
\left\{
N=N_\varepsilon,
\ Y-a_\varepsilon'
\in\delta_\varepsilon B
\right\}
=\mu_\varepsilon^{-1/2}
\left(c_{H,a}+o(1)\right),
\]
where \(c_{H,a}\) is continuous and bounded above and below by positive
constants.  More generally, if the continuous shell is on the Gaussian scale,
its limiting Gaussian mass replaces one; for narrower shells the corresponding
mixed local-density formula holds under the usual boundary-smoothing scale.
In every case used for microcanonical conditioning, the shell logarithm is
\(o(\mu_\varepsilon)\).
\end{theorem}

\begin{proof}
Use Fourier inversion in the exact number coordinate and smooth upper and
lower approximations to \(1_{\delta_\varepsilon B}\).  In the central region,
expand the pressure at the exact finite saddle.  The lattice coefficient gives
the Schur-complement one-dimensional factor
\((2\pi\mu_\varepsilon\sigma_{N|Y}^2)^{-1/2}\).  The continuous shell contains
an increasing number of standard deviations and its Gaussian mass tends to
one.  Boundary smoothing contributes
\(o(\mu_\varepsilon^{-1/2})\) because the boundary is piecewise \(C^2\) and
the smoothing width can be chosen between
\(\mu_\varepsilon^{-1/2}\) and \(\delta_\varepsilon\).  The three remaining
Fourier regions are negligible by
Theorem~\ref{thm:r6-b1-characteristic}.  Uniform covariance bounds give the
positive bounds on \(c_{H,a}\).
\end{proof}

\subsection{Source-dependent conditioning}

Let \(\mathcal S_\varepsilon(a_\varepsilon)\) denote the exact number and
continuous shell above.  Define
\[
Q_\varepsilon^{\rm mc,a}(H)=
\frac1{\mu_\varepsilon}\log
\mathbb E_\varepsilon^{\rm gc}
\left[e^{\mu_\varepsilon\langle H,X_\varepsilon\rangle}
\mid\mathcal S_\varepsilon(a_\varepsilon)\right].
\]

\begin{theorem}[Microcanonical pressure and covariance]
\label{thm:r6-b1-transfer}
Locally uniformly in \(H\),
\[
Q_\varepsilon^{\rm mc,a}(H)\longrightarrow
Q(H,\lambda_H)-\lambda_H\cdot a
-Q(0,\lambda_0)+\lambda_0\cdot a.
\]
Equivalently this is the difference of the two constrained infima in
\(\lambda\).  Its Hessian is
\[
Q_{HH}-Q_{H\lambda}Q_{\lambda\lambda}^{-1}Q_{\lambda H},
\]
the exact constraint Schur complement.  For a time-zero constraint source
\(H_\eta\), the pressure equals \(\eta\cdot a\).
\end{theorem}

\begin{proof}
Change measure in the numerator by the exact finite-volume tilt
\((H,\lambda_{H,\varepsilon})\).  On the shell, the multiplier term is
\(-\mu_\varepsilon\lambda_{H,\varepsilon}\cdot a_\varepsilon+o(\mu_\varepsilon)\),
and Theorem~\ref{thm:r6-b1-coefficient} makes the remaining shell probability
subexponential.  Repeat with \(H=0\) for the denominator and pass to the
uniform saddle limit.  Differentiating the finite saddle equation and then the
envelope gives the Schur complement.  For a time-zero source,
\(Q(H_\eta,\lambda)=Q(0,\lambda+\eta)\), so the infimum translates exactly.
\end{proof}

\subsection{The prepared initial LDP}

Use the separating class of bounded time-zero cylinder sources together with a
small quadratic velocity chart.  Maxwellian exponential tightness gives
compact empirical-measure sublevels.

\begin{theorem}[Round-six B1 closure]
\label{thm:r6-b1-main}
The primitive hard-sphere shell is conditioned at a unique exact finite-volume
source-dependent saddle.  The singleton compound-Poisson exponent supplies a
uniform mixed lattice/nonlattice local limit, the shell coefficient has only
subexponential cost, and the limiting pressure, initial constrained rate, and
Gaussian covariance are the source-dependent infimum and Schur-complement
formulas above.
\end{theorem}

\begin{proof}
The mixed shell theorem and pressure transfer are
Theorems~\ref{thm:r6-b1-coefficient} and
\ref{thm:r6-b1-transfer}.  Exponential tightness and the conditional Laplace
principle on the separating time-zero source class give the good initial LDP;
its rate is the grand-canonical initial rate restricted to the constraint
surface and normalized by its constrained minimum.  The covariance statement
was proved in Theorem~\ref{thm:r6-b1-transfer}.
\end{proof}
